\documentclass[11pt,a4paper]{article}

% --- Encodage / langue ---
\usepackage[utf8]{inputenc}
\usepackage[T1]{fontenc}
\usepackage[french]{babel}
\usepackage{lmodern}

% --- Mise en page compacte (tient en 2 pages hors illustrations) ---
\usepackage[a4paper,margin=1.9cm]{geometry}
\usepackage{microtype}
\usepackage{enumitem}
\setlist{nosep,leftmargin=1.2em}
\usepackage{amsmath}
\usepackage{booktabs}
\usepackage{titlesec}
\titlespacing*{\section}{0pt}{6pt}{2pt}
\titlespacing*{\subsection}{0pt}{4pt}{1pt}
\titleformat{\section}{\normalfont\large\bfseries}{\thesection.}{0.5em}{}
\titleformat{\subsection}{\normalfont\bfseries}{}{0em}{}
\setlength{\parindent}{0pt}
\setlength{\parskip}{2pt}

\usepackage[backend=biber,style=numeric-comp,sorting=none,maxbibnames=3]{biblatex}
\addbibresource{references.bib}
\renewcommand*{\bibfont}{\footnotesize}

\usepackage{hyperref}
\hypersetup{hidelinks}

\title{\vspace{-1.2cm}\textbf{Livrable 3 --- Prototype Python : description et évaluation\\[2pt]
\large Analyse automatique de partitions musicales (OMR)}}

\author{\normalsize Thibault Bourdoiseau \& Léo Li --- SAE TIFO 1 / TIFO 2 / MLRF}
\date{\vspace{-8pt}\small \today}

\begin{document}
\maketitle
\vspace{-10pt}

\section{Périmètre retenu}
Le prototype transforme l'image d'une partition \textbf{monophonique} en une séquence de notes
(hauteur~+~durée) exportée en MIDI, avec les seules notions TIFO~1, TIFO~2 et MLRF~\cite{rebelo2012,calvo2020}
--- \textbf{sans réseau de neurones}. L'état de l'art et le livrable~2 balisaient large (imprimé \emph{et}
manuscrit, détection \emph{et} décodage). Nous avons fait le choix assumé de livrer une \textbf{chaîne
complète et robuste sur un cas} --- l'imprimé et le déformé de \textbf{CameraPrIMuS} ($87\,678$ incipits, image
$\rightarrow$ MIDI) --- plutôt qu'un prototype fragile étalé sur trois bases. Le manuscrit (MUSCIMA++~\cite{hajic2017})
et le F1 de détection sur pages polyphoniques (DeepScoresV2~\cite{tuggener2020}) sont identifiés comme
extensions hors périmètre du prototype (cf.~§\ref{sec:limites}). Code : NumPy + Pillow uniquement, pour un
portage C++ direct au WP3.

\section{Chaîne de traitement retenue}
L'image suit une chaîne classique~\cite{rebelo2012}~; chaque étape ne mobilise que des notions de cours
(traçabilité complète dans \texttt{CORRESPONDANCE\_COURS.md}).
\begin{enumerate}
\item \textbf{Prétraitement} (\texttt{preprocess.py}, TIFO~1). Niveaux de gris $\rightarrow$ \textbf{filtre médian}
(bruit) $\rightarrow$ binarisation par le \textbf{critère d'Otsu}~\cite{otsu1979} (imprimé) ou \textbf{seuil
adaptatif} par image intégrale (photos à éclairage non uniforme).
\item \textbf{Portées} (\texttt{staff.py}, TIFO~1/2). Épaisseur de ligne et interligne estimés par
\textbf{longueurs de runs}~; les $5$ lignes localisées par \textbf{projection horizontale}. Suppression par
\textbf{ouverture morphologique} (élément structurant vertical) \emph{unie} à un mécanisme de « pont » qui
conserve les pixels là où un symbole traverse la ligne (\texttt{remove\_staff\_lines\_combined}) --- ce qui
débloque les têtes creuses et les altérations.
\item \textbf{Segmentation} (\texttt{segment.py}, TIFO~1/2). \textbf{Composantes connexes} 8-connexité
(union-find à deux passes, NumPy pur) $\rightarrow$ une composante par symbole.
\item \textbf{Détection des têtes} (\texttt{heads.py}, TIFO~2). Bouchage des petits trous (reconstruction
géodésique) $\rightarrow$ \textbf{ouverture par un disque} (granulométrie) $\rightarrow$ filtrage par taille,
allongement et remplissage. Détecte les têtes y compris dans les groupes ligaturés.
\item \textbf{Hauteurs} (\texttt{decode.py}). Géométrie de la portée (ligne du bas + demi-interligne)
$\rightarrow$ position verticale de la tête $\rightarrow$ lettre + octave via la référence de la clef.
\item \textbf{Durées} (\texttt{duration\_knn.py}, MLRF~L1/L3). \textbf{Descripteur de forme} (histogrammes de
projection + moments de Hu + géométrie) de la région tête+hampe+ligature, classé au \textbf{$k$ plus proches
voisins} (ronde/blanche/noire/croche/\dots). Repli géométrique conservé si aucun modèle n'est fourni.
\item \textbf{Export MIDI} (\texttt{midi.py}). Écriture d'un fichier Standard MIDI pur (hauteur $\rightarrow$
numéro MIDI, durée $\rightarrow$ tics), sans dépendance externe.
\end{enumerate}
La \textbf{clef} est reconnue par descripteurs + $k$-NN ($99{,}2\,\%$ imprimé, $91{,}7\,\%$ distordu)~;
l'\textbf{armure} et les \textbf{silences} disposent de méthodes conformes (théorie + $k$-NN de forme) mais leur
détection automatique reste \textbf{désactivée / opt-in} (cf.~§\ref{sec:limites}).

\subsection{Décisions et écarts assumés vs état de l'art}
\begin{itemize}
\item \textbf{Suppression de portée par ES vertical} (l'état de l'art décrit un ES \emph{horizontal}, qui abîme
les ligatures --- vérifié)~: reste une ouverture morphologique (TIFO2~§5.3).
\item \textbf{Durées par $k$-NN} (annoncées géométriques au §5)~: extension de la méthode $k$-NN déjà prévue
au §4, la géométrie seule s'étant révélée peu fiable à cette résolution ($\sim$3~px de ligature). \emph{Gain majeur.}
\item \textbf{Détection des têtes}~: les deux options du §4 testées (morphologie vs corrélation de gabarit)~;
la morphologie gagne nettement (ratio têtes/notes $1{,}07$ vs $2{,}1$) et est retenue.
\item \textbf{Portées inclinées/courbes}~: \textbf{Hough} et \textbf{chemins stables}~\cite{cardoso2009}
implémentés et validés sur le déformé (Hough trouve les $5$ lignes dans $99\,\%$ des cas vs $61\,\%$ pour la
projection)~; la projection reste le défaut sur imprimé droit (plus simple, suffisante).
\end{itemize}

\section{Évaluation}
Protocole~: critères PA / RA / NA du livrable~2 (alignement par distance d'édition de
Levenshtein~\cite{levenshtein1966}), sur CameraPrIMuS, $n=800$ incipits, \textbf{split train/test disjoint}
$70/30$ (le classifieur de durées n'est entraîné que sur \emph{train}, l'évaluation ne porte que sur
\emph{test} $=240$). Sous-ensembles~: \emph{simple} (sans armure ni notes creuses) et \emph{général} (tout).
La ligne \emph{armure oracle} ($n=400$) fournit la vérité-terrain de l'armure au décodage~: elle borne le gain
théorique d'une lecture parfaite des altérations. On situe l'algorithme retenu (\textbf{V2}, durées par
$k$-NN) par rapport à la \textbf{baseline V1}~: le décodage des durées par \textbf{règles géométriques}
(tête pleine/vide, hampe, comptage de ligatures et points), c'est-à-dire la méthode décrite dans notre état de
l'art (§5). V2 n'est donc pas une méthode étrangère mais une \emph{amélioration mesurée} de cette baseline.

\begin{center}\small
\begin{tabular}{llccccc}
\toprule
Sous-ensemble & Version & $N$ & PA stricte & PA diat. & RA & NA \\
\midrule
simple  & V1 (géom., état de l'art §5) & 61  & 75{,}5 & 77{,}1 & 40{,}9 & 38{,}5 \\
simple  & \textbf{V2 (k-NN, retenu)} & 61  & \textbf{83{,}7} & \textbf{85{,}7} & \textbf{77{,}4} & \textbf{74{,}3} \\
général & V1 (géom., état de l'art §5) & 240 & 60{,}4 & 76{,}7 & 39{,}4 & 29{,}3 \\
général & \textbf{V2 (k-NN, retenu)} & 240 & \textbf{66{,}1} & \textbf{83{,}7} & \textbf{71{,}4} & \textbf{56{,}2} \\
général & V2 + armure oracle & 120 & 82{,}1 & 84{,}3 & 69{,}4 & \textbf{66{,}2} \\
\bottomrule
\end{tabular}
\end{center}

\textbf{Lecture.} Le passage V1~$\rightarrow$~V2 (classifieur de durées $k$-NN) est le gain majeur~: sur les
notes bien alignées, la précision de durée passe de \textbf{$39\,\%$ à $79\,\%$}, ce qui double
quasiment la NA générale. La PA monte aussi mécaniquement (des durées correctes améliorent l'alignement par
distance d'édition). Sur le sous-ensemble \emph{simple}, le décodage est presque parfait (PA diat.\ $85{,}7$).
L'écart PA stricte / diatonique et la ligne \emph{oracle} chiffrent le \textbf{plafond restant}~: la lecture
des altérations (armure). Tous les essais intermédiaires, gardés comme abandonnés, sont journalisés avec leurs
scores dans \texttt{SUIVI.md}.

\textbf{Cohérence avec l'état de l'art.} En OMR \emph{classique} (sans réseau), une exactitude note-level de
$50$--$80\,\%$ sur données réelles est attendue~\cite{rebelo2012,calvo2020}~; seules les approches réseau
dépassent $\sim$$90\,\%$ sur PrIMuS. Nos chiffres (NA simple $74$, général $56$, PA diat.\ $84$) se situent
donc dans la fourchette haute du classique.

\section{Limites et reproductibilité}\label{sec:limites}
\textbf{Limites (honnêtes).} (i)~\emph{Armure / altérations} = plafond n\textsuperscript{o}1~: théorie et
intégration faites (l'oracle le chiffre), mais aucune détection automatique n'atteint une précision
nette-positive (la suppression de portée fragmente les traits fins $\sim$3~px) $\rightarrow$ laissée OFF pour ne
pas dégrader la NA. (ii)~\emph{Silences}~: méthode implémentée (candidats + $k$-NN de forme) mais sur-détection
$\rightarrow$ opt-in. (iii)~\emph{Manuscrit (MUSCIMA++)} et \emph{F1 sur pages polyphoniques (DeepScoresV2)}~:
hors périmètre --- la chaîne de décodage suppose une portée unique monophonique~; le F1 est rapporté par un
proxy (ratio têtes/notes $\approx 1{,}07$) faute d'annotations de boîtes dans PrIMuS. Ce sont les vraies limites
d'un prototype sans réseau.

\textbf{Reproductibilité.} Environnement géré par \texttt{uv}. Tests unitaires~: \texttt{uv run pytest -q}
($34$ tests). Tableau ci-dessus~: \texttt{uv run python scripts/benchmark.py -{}-n 800} (V2),
\texttt{-{}-geom} (V1), \texttt{-{}-key-oracle} (oracle). Démo image~$\rightarrow$~MIDI~:
\texttt{uv run python -m omr.cli -{}-id <incipit>}.

\printbibliography[heading=bibnumbered,title={Références}]

\end{document}
